% Bagian: sets-functions-relations
% Bab: size-of-sets
% Subbagian: pairing

\documentclass[../../../include/open-logic-section]{subfiles}

\begin{document}

\olfileid[id]{sfr}{siz}{pai}
\olsection{Fungsi Pemasangan dan Kode}

\begin{explain}
Metode zig-zag Cantor membuat keterhitungan $\Nat^n$ tampak jelas secara
visual. Namun, mari kita pusatkan perhatian pada larik yang menggambarkan
$\Nat^2$. Dengan mengikuti garis zig-zag pada larik dan menghitung posisinya,
kita dapat memeriksa bahwa $\tuple{1,2}$ dipasangkan dengan bilangan~$7$.
Meskipun demikian, akan lebih baik jika kita dapat menghitungnya secara lebih
langsung. Dengan kata lain, akan berguna jika kita mempunyai \emph{invers}
dari enumerasi zig-zag tersebut, $g\colon \Nat^2 \to \Nat$, sedemikian sehingga
\[
g(\tuple{0,0}) = 0, \;
g(\tuple{0,1}) = 1, \;
g(\tuple{1,0}) = 2, \; \dots,
g(\tuple{1,2}) = 7, \; \dots
\]
Fungsi ini akan memungkinkan kita menghitung tepat di posisi mana
$\tuple{n, m}$ muncul dalam enumerasi kita.

Sebenarnya, kita dapat mendefinisikan $g$ secara langsung dengan membuat dua
pengamatan. Pertama: jika baris ke-$n$ dan kolom ke-$m$ memuat nilai~$v$,
maka baris ke-$(n+1)$ dan kolom ke-$(m-1)$ memuat nilai $v + 1$. Kedua: baris
pertama enumerasi kita terdiri atas bilangan segitiga, dimulai dengan $0$,
$1$, $3$, $6$, dan seterusnya. Bilangan segitiga ke-$k$ adalah jumlah bilangan
bulat positif hingga dan termasuk $k$, yang dapat dihitung sebagai
$k(k+1)/2$. Dengan
menggabungkan kedua pengamatan tersebut, perhatikan fungsi berikut:
\[
  g(n,m) = \frac{(n+m+1)(n+m)}{2} + n
\]
Kita sering menulis $g(n, m)$ saja alih-alih $g(\tuple{n, m})$ karena lebih
mudah dibaca. Rumus ini meminta kita terlebih dahulu menentukan bilangan
segitiga urutan $(n+m)^\text{ke}$, lalu menambahkan $n$. Rumus tersebut mengisi larik tepat
seperti yang kita inginkan. Jadi, secara khusus, pasangan $\tuple{1, 2}$
dipetakan ke $\frac{4 \times 3}{2} + 1 = 7$.

Fungsi $g$ ini adalah \emph{invers} dari enumerasi suatu himpunan pasangan.
Fungsi semacam itu disebut \emph{fungsi pemasangan}.
\end{explain}

\begin{defn}[Fungsi pemasangan]
  Suatu fungsi $f\colon A \times B \to \Nat$ adalah \emph{fungsi pemasangan}
  aritmetis jika $f$ injektif. Kita juga mengatakan bahwa $f$
  \emph{mengodekan} $A \times B$, dan bahwa $f(x,y)$ adalah \emph{kode} bagi
  $\tuple{x,y}$.
\end{defn}

\begin{explain}
Kita dapat menggunakan fungsi pemasangan untuk mengodekan, misalnya, pasangan
bilangan asli; dengan kata lain, kita dapat merepresentasikan setiap
\emph{pasangan} anggota-anggota dengan sebuah bilangan \emph{tunggal}. Dengan
menggunakan invers fungsi pemasangan, kita dapat \emph{mendekodekan} bilangan
tersebut, yakni menentukan pasangan yang direpresentasikannya.
\end{explain}

\begin{prob}
Berikan suatu enumerasi untuk himpunan semua bilangan rasional taknegatif.
\end{prob}

\begin{prob}
Tunjukkan bahwa $\Rat$ adalah !!{enumerable}. Ingat bahwa setiap bilangan
rasional dapat ditulis sebagai pecahan $z/m$ dengan $z \in \Int$, $m \in \Nat^+$.
\end{prob}

\begin{prob}
Definisikan suatu enumerasi untuk $\Bin^*$.
\end{prob}

\begin{prob}
Ingat dari mata kuliah pengantar logika bahwa setiap tabel kebenaran yang
mungkin menyatakan suatu fungsi kebenaran. Dengan kata lain, fungsi-fungsi
kebenaran adalah semua fungsi dari $\Bin^k \to \Bin$ untuk suatu~$k$.
Buktikan bahwa himpunan semua fungsi kebenaran adalah terhitung.
\end{prob}

\begin{prob}
Tunjukkan bahwa himpunan semua himpunan bagian hingga dari suatu himpunan
tak hingga yang !!{enumerable} adalah !!{enumerable}.
\end{prob}

\begin{prob}
Suatu himpunan bagian dari $\Nat$ disebut \emph{kofinit} jika dan hanya jika
himpunan itu merupakan komplemen dalam $\Nat$ dari suatu himpunan hingga;
yakni, $A \subseteq \Nat$ kofinit jika dan hanya jika $\Nat\setminus A$
hingga. Misalkan $I$ adalah himpunan yang !!{element}snya tepat semua himpunan
bagian hingga dan kofinit dari $\Nat$. Tunjukkan bahwa $I$ adalah
!!{enumerable}.
\end{prob}

\begin{prob}
Tunjukkan bahwa gabungan !!{enumerable} dari himpunan-himpunan yang !!{enumerable}
adalah !!{enumerable}. Artinya, jika $A_1$, $A_2$, \dots{} adalah himpunan dan
setiap $A_i$ adalah !!{enumerable}, maka gabungan $\bigcup_{i=1}^\infty A_i$
dari semuanya juga !!{enumerable}. [Catatan: ini sulit!]
\end{prob}

\begin{prob}
Misalkan $f \colon A \times B \to \Nat$ adalah sebarang fungsi pemasangan.
Tunjukkan bahwa invers parsial dari $f$ pada daerah hasilnya menentukan suatu
enumerasi untuk $A \times B$.
\end{prob}

\begin{prob}
Tentukan suatu fungsi yang mengodekan $\Nat^3$.
\end{prob}

\end{document}
